\SGASubsection{7.10}{}
\SGARefTarget{sga3:xxvi:subsec:7:10}

The relative positions of two minimal parabolic subgroups can be
described as follows.  We use the notation \(t_2(P,P_1)\) of 4.5.2.

\begin{enumerate}[label=\textup{(\arabic*)}]
\item Let \(P,P_1,P',P'_1\) be four minimal parabolic subgroups of
\(G\).  Then
\[
 t_2(P,P_1)=t_2(P',P'_1),
\]
that is, the pairs \((P,P_1)\) and \((P',P'_1)\) are fpqc-locally
conjugate, if and only if there is a \(g\in G(S)\) such that
\[
 \Int(g)P=P',\qquad \Int(g)P_1=P'_1.
\]

\item Fix a minimal parabolic subgroup \(P\) of \(G\) with Levi
subgroup \(L\), and let \(T\) be a maximal torus of \(L\).  Consider
the scheme
\[
 \Par_{t_{\min}}(G;P)
\]
of minimal parabolic subgroups of \(G\) in standard position
relative to \(P\).  There is a morphism \textup{(cf.\ \SGARefLink{sga3:xxvi:paragraph:4-5-5}{4.5.5})}
\[
 f:\Par_{t_{\min}}(G;P)
   \longrightarrow W_P(T)\backslash W_G(T)/W_P(T)
\]
whose fibers are the orbits of \(P\) in
\(\Par_{t_{\min}}(G;P)\).  By \textup{(1)}, \(f\) therefore induces
a monomorphism
\[
 P(S)\backslash\Par_{t_{\min}}(G;P)(S)
 \lhook\joinrel\longrightarrow
 \bigl(W_P(T)\backslash W_G(T)/W_P(T)\bigr)(S).
\]
Its image is identified with \(W\).  This is the Bruhat theorem:
each orbit of \(P(S)\) in \(\Par_{t_{\min}}(G;P)(S)\) contains
exactly one parabolic subgroup of \(G\) with Levi subgroup \(L\),
that is, one of the form
\[
 \Int(n)P,\qquad n\in\Norm_G(L)(S).
\]

\item Equivalently, let \(E(S)\) be the set of elements \(g\in G(S)\)
such that \(\Int(g)P\) and \(P\) are in mutual standard position.
Then \(E(S)\) is partitioned into double cosets modulo \(P(S)\),
indexed by \(W\):
\[
 E(S)=P(S)\,W\,P(S).
\]
Putting \(U=\rad^u(P)\), one may also write
\[
 E(S)=U(S)\cdot\Norm_G(L)(S)\cdot U(S).
\]
This exhibits a partition of \(E(S)\) into double cosets modulo
\(U(S)\), indexed by \(\Norm_G(L)(S)\).

\item If \(S\) is the spectrum of a field, then \(E(S)=G(S)\), and
one recovers \textup{[BT65]}, 5.15.
\end{enumerate}

\begin{sgastatement}{Counter\SGARefLink{sga3:xxvi:counterexamples:7:11}{examples 7.11}}
Let
\[
 S=\operatorname{Spec}(\mathbf Z/4\mathbf Z),
 \qquad G=\operatorname{SL}_{2,S}.
\]
Let \(B\) be the usual Borel subgroup consisting of the matrices
\[
 \begin{pmatrix}a&b\\c&d\end{pmatrix}
 \qquad\text{with }c=0,
\]
and put
\[
 g=\begin{pmatrix}-1&1\\2&1\end{pmatrix}\in G(S),
 \qquad B'=\Int(g)B.
\]
Then
\[
 B(S)=B'(S),
\]
whereas \(B\cap B'\) contains no maximal torus.
\footnote{Indeed, \(B'\) is the subgroup of \(G\) defined by
\(c=2(a+b)\).  Thus \(B\cap B'\) is not flat over \(S\), and hence
contains no maximal torus, by \SGARefLink{sga3:xxvi:proposition:4-5-1}{4.5.1}.}

This shows both that two distinct minimal parabolic subgroups may
have the same group of sections and that, in general, one cannot
decide from the groups \(P(S)\) and \(P'(S)\) alone whether two
minimal parabolic subgroups \(P,P'\) are in standard position.  In
particular, the subset \(E(S)\) of \(G(S)\) does not appear to be
definable solely from the configuration
\[
 \{G(S),P(S),\Norm_G(L)(S)\}.
\]
In the example above, this subset is defined by \(c\ne2\).
\footnote{More generally, for every \(S\)-scheme \(S'\), the set
\(E(S')\) is defined by the condition that \(c\) is either zero or
invertible.}
\end{sgastatement}
\SGARefTarget{sga3:xxvi:counterexamples:7:11}
